\documentclass[10pt,a4paper]{amsart}
\usepackage[margin=19mm]{geometry}
\usepackage{amsmath,amssymb,mathtools}
\usepackage[scr=rsfs,cal=euler]{mathalfa}
\usepackage{xfrac}
\usepackage[unicode,hidelinks,bookmarks=false]{hyperref}
\newcommand{\ie}{i.e.\ }
\newcommand{\cf}{cf.\ }
\newcommand{\resp}{respectively\ }
\newcommand{\Subsection}{\subsection}
\providecommand{\nobreakdash}{\nobreak\hbox{-}}
\setlength{\emergencystretch}{2em}
\multlinegap0pt
\usepackage{bm}
\usepackage{bbm}
\usepackage{dsfont}
\usepackage{xspace}
\usepackage{cancel}
\usepackage{mathtools}
\usepackage{amsthm}
\makeatletter
\@ifundefined{theorem}{\newtheorem{theorem}{Theorem}}{}
\@ifundefined{lemma}{\newtheorem{lemma}[theorem]{Lemma}}{}
\@ifundefined{proposition}{\newtheorem{proposition}[theorem]{Proposition}}{}
\makeatother
\newcommand{\hra}{\hookrightarrow}
\newcommand{\ol}{\overline}
\renewcommand{\d}{{\mbox{d}}}
\renewcommand{\subset}{\subseteq}
\title[A note on the cohomology of pullback Lie algebroids]{A note on the cohomology of pullback Lie algebroids}
\author{Eckhard Meinrenken}
\date{Source: arXiv:2607.23859v1 (CC BY 4.0)}
\begin{document}
\maketitle

\noindent\emph{Source and licence.} A note on the cohomology of pullback Lie algebroids. Version arXiv:2607.23859v1 (CC BY 4.0), distributed under the Creative Commons Attribution 4.0 International licence, as recorded in the arXiv licence metadata for this exact version and shown on the abstract page https://arxiv.org/abs/2607.23859v1. Full rights notice: licenses/arxiv-2607.23859-CC-BY-4.0.txt.\par\smallskip

\noindent\textit{Editorial scope.} Counted unit: the Lemma whose source label is \texttt{lem:3.6} in the article (source0.tex lines 460--465, source label \texttt{lem:3.6}), delivered together with the article's own complete original proof of it (source0.tex lines 468--501, one \texttt{proof} environment of 34 lines). No mathematics is written, repaired or re-derived: statement and proof are the article's own text line for line. Required context retained uncounted: prop:product (lines 338--342). Every definition, display and label of those lines is retained unchanged. Omitted from this package and not redistributed: 1--337, 343--459, 466--467, 502--943. Nothing inside a retained excerpt is omitted or altered: every retained body is delivered exactly as the article's lines are written, so no omission receipt of any kind accompanies this package. Citations: the retained excerpts carry no citation command at all, so no bibliography entry of the article is transcribed and no reference is redistributed. Reference adaptation: none was needed and none was made; every reference inside the delivered unit resolves inside it, so no unresolved reference or citation is produced. Printed slips kept literal: no printed word, symbol, hypothesis or number of the counted statement or of its proof was altered. Layout and class: the article's own class is \texttt{amsart} with packages including amsmath, amscd, amssymb, amsfonts, amsthm, hyperref, tcolorbox, tikz-cd, xy, mathabx, graphicx, MnSymbol; the wrapper is the lane's pinned \texttt{amsart} editorial wrapper at 10pt on A4, which restates the theorem-like environments the retained text needs and adds the article's own macro definitions that the retained text needs (\texttt{\textbackslash d}, \texttt{\textbackslash hra}, \texttt{\textbackslash ol}, \texttt{\textbackslash subset}), with extra wrapper packages (bm, bbm, dsfont, xspace, cancel, mathtools). The delivered numbering is the unit's own. Assets: no retained span contains a figure environment, a graphics-inclusion command, a PDF-page inclusion, a table or a TikZ picture; no image file is copied into this package, the package declares no asset dependency and no image file is redistributed with it. Licence: version arXiv:2607.23859v1 of this article is distributed under the Creative Commons Attribution 4.0 International licence, as recorded in the arXiv licence metadata for this exact version and shown on the abstract page https://arxiv.org/abs/2607.23859v1. A full rights notice is delivered at licenses/arxiv-2607.23859-CC-BY-4.0.txt. Characters: every retained excerpt is pure ASCII, so the delivered engine, XeLaTeX with its default Latin Modern text font, reports no missing character.
\begin{proposition}[Products] \label{prop:product}
	Suppose $M=Q\times N$ 	is a product of two manifolds, where $Q$ is of finite type, and let $\pi\colon M\to N$ be the projection to the second factor. 
	Then 
	\[ H_\pi^\bullet(M)=H^\bullet(Q)\otimes C^\infty(N).\]
\end{proposition}

\begin{lemma}\label{lem:3.6}
	Suppose that for some $k$, we have that 
	$H^k(\pi^{-1}(x))=0$ for all $x\in N$. Given an exhaustion as in the previous lemma, there exists, for all $a$, some $b>a$ such that the restriction map  
	\[ H^k_\pi(U_{b})\to H^k_\pi(U_a)\]
	is zero. 
\end{lemma}

\begin{proof}\phantom{.}
%\begin{enumerate}\renewcommand{\theenumi}{\arabic{enumi}}	\renewcommand{\labelenumi}{\theenumi.}
	%\item 	
	Given $x\in N$, fix  $c>a+1$ so that 	
	\[ H^k(U_{c}\cap \pi^{-1}(x))\to H^k(U_{a+1}\cap \pi^{-1}(x))\]
	is the zero map. 
	%\item 
	Choose a product open neighborhood $\Phi((U_c\cap \pi^{-1}(x))\times V)$ as in 
	the previous lemma. 
	%\item 
	Shrinking $V$ if necessary, the product neighborhood  $\Phi((U_{a+1}\cap \pi^{-1}(x))\times V)$ contains 
	$\ol{U}_{a}\cap \pi^{-1}(V)$, and taking $b$ sufficiently large the neighborhood $\Phi((U_c\cap \pi^{-1}(x))\times V)$ 
	is contained in $U_b$. This gives inclusions
	\[ U_{a}\cap \pi^{-1}(V)\hra  \Phi((U_{a+1}\cap \pi^{-1}(x))\times V)\hra \Phi((U_c\cap \pi^{-1}(x))\times V)
	\hra U_b\cap \pi^{-1}(V)
	.\]
   %\item 
  The map 
   \[ H^k_\pi (\Phi((U_c\cap \pi^{-1}(x))\times V))\to  H^k_\pi (\Phi((U_{a+1}\cap \pi^{-1}(x))\times V)),\]
 induced by the middle inclusion is the zero map, since 
 \[ H^k_\pi (\Phi((U\cap \pi^{-1}(x))\times V))=H^k(U\cap \pi^{-1}(x))\otimes C^\infty(V)\]
 by Proposition \ref{prop:product}. It follows that the map 
   %
   \[ H^k_\pi(U_{b}\cap \pi^{-1}(V))\to H^k_\pi(U_a\cap \pi^{-1}(V))\]
 is the zero map.  The base manifold $N$ may be covered by open subsets $V$ of this sort. Since 
   $U_a$ has compact closure, there is a finite collection $\{V_i\}$ of such open subsets, such that 
   $\{\pi^{-1}(V_i)\}$ cover $\ol{U}_a$. Choose $b_i$ such that the maps 
     \[ H^k_\pi(U_{b_i}\cap \pi^{-1}(V_i))\to H^k_\pi(U_a\cap \pi^{-1}(V_i))\]
  are all zero, and let $b$ be their maximum. Given $\alpha\in \Omega^k_\pi(U_b)$ with $\d_\pi\alpha=0$, 
  the restriction to $U_{a}\cap \pi^{-1}(V_i)$ admits a primitive $\beta_i$. 
  Let  $\{\chi_i\}$ be a partition of unity  subordinate to $\{V_i\}$ (regarded as a cover of $\bigcup V_i\subset N$) 
    and put $\beta=\sum_i (\pi^*\chi_i)\beta_i$. 
  Then $\d_\pi\beta=\alpha$ on $U_a$. 
\end{proof}

\end{document}
